This study examines how Pakistani American immigrant families construct and negotiate identity through retrospective storytelling across generations. This study integrates Communication Theory of Identity (CTI) and Communicated Narrative Sense-Making (CNSM) to demonstrate that storytelling is the communication process through which layered identities are actively produced and negotiated. 

Using a qualitative, interpretivist Art Based Research (ABR) approach, this study employs episodic narrative in-depth interviews with 24 first-generation parents and 1.5/second-generation adult children from Pakistani American families. poetic inquiry as a themeatic analysis tool. 

Findings revealed that retrospective storytelling both reveals and mediates identity gaps, enabling individuals to navigate the tensions between who they are, who they are expected to be, and how they are perceived across contexts. The study introduces Communal Retrospective Storytelling and its function as shared decision-making infrastructure through which immigrant families collectively interpret past experiences of Pakistani community members in their surroundings to manage present uncertainty and anticipate future outcomes of their children.  
